\section{Constraint-Solver Reading Layer}

\subsection{Purpose}

The project has reached the point where speed depends on automatic rejection. This section adds a constraint-solver layer that converts anchor dossiers into ranked reading skeletons. The solver is deliberately pre-phonetic: it can accept a structural reading such as \texttt{QUALIFIER/TITLE + ENTITY STEM}, but it cannot assign a sound value, language-family value, or lexical translation.

This is the safest fast path toward decipherment. It lets us scale from a few promising anchors to many hypotheses while keeping each proposal falsifiable.

\subsection{Method}

The script \texttt{scripts/constraint-solver-reading-layer.py} combines:

\begin{itemize}
  \item anchor dossiers and component roles;
  \item occurrence-level context rows;
  \item minimal-pair and component-alternation evidence;
  \item lexical reading gates;
  \item structural frames and semantic frame labels.
\end{itemize}

Each candidate is tested against seven constraints: component order, context support, minimal-pair support, frame stability, lexical-gate compliance, phonetic restraint, and functional-role availability. The output is a ranked list of skeleton readings, not deciphered sentences.

\subsection{Solver Result}

\begin{table}[htbp]
\centering
\begin{tabular}{lr}
\toprule
\textbf{Metric} & \textbf{Value} \\
\midrule
Candidate anchor skeletons & 11 \\
Accepted/promising skeletons & 6 \\
Exploratory skeletons & 5 \\
Hard-constraint holds & 0 \\
Clause reconstructions exported & 59 \\
Phonetic readings assigned & 0 \\
\bottomrule
\end{tabular}
\caption{Constraint-solver summary.}
\label{tab:constraint-solver-summary}
\end{table}

The strongest accepted skeleton is:

\begin{center}
\texttt{240-482 = QUALIFIER/TITLE PREFIX + ENTITY STEM}
\end{center}

The most important supporting chain is:

\begin{center}
\texttt{176 = ENTITY STEM}
\qquad
\texttt{240-176 = QUALIFIER/TITLE PREFIX + ENTITY STEM}
\end{center}

This is not a translation. It is a controlled structural claim: \texttt{240} behaves like a reusable qualifier/title/classifier-like component, while \texttt{176} behaves like a reusable root-like entity component.

\subsection{Highest-Ranked Reading Skeletons}

\begin{table}[htbp]
\centering
\footnotesize
\begin{tabular}{lp{0.48\textwidth}lr}
\toprule
\textbf{Unit} & \textbf{Skeleton} & \textbf{Tier} & \textbf{Score} \\
\midrule
\texttt{240-482} & \texttt{QUALIFIER\_TITLE\_PREFIX + ENTITY\_STEM} & A & 0.966 \\
\texttt{240-740-090} & \texttt{ADMINISTRATIVE\_OBJECT\_FORMULA} & A & 0.951 \\
\texttt{032} & \texttt{TERMINAL\_TITLE\_OR\_OBJECT\_ENDING} & A & 0.882 \\
\texttt{176} & \texttt{ENTITY\_STEM} & B & 0.849 \\
\texttt{240-176} & \texttt{QUALIFIER\_TITLE\_PREFIX + ENTITY\_STEM} & B & 0.783 \\
\texttt{590-390} & \texttt{QUALIFIER\_TITLE\_PREFIX + FINAL\_QUALIFIER} & B & 0.757 \\
\bottomrule
\end{tabular}
\caption{Highest-ranked pre-phonetic reading skeletons.}
\label{tab:constraint-solver-top-skeletons}
\end{table}

\subsection{Approximate Distance to Full Decipherment}

Full decipherment remains far. The project is no longer at zero: we have corpus control, structural frames, abstract variables, semantic context scores, anchor dossiers, and now constrained reading skeletons. But a strict decipherment requires validated sound values and a language identification, and those remain mostly unearned.

As an approximate evidence-maturity estimate:

\begin{table}[htbp]
\centering
\begin{tabular}{lr}
\toprule
\textbf{Subproblem} & \textbf{Maturity} \\
\midrule
Corpus and provenance control & 0.70 \\
Structural grammar & 0.48 \\
Morpheme segmentation & 0.27 \\
Semantic anchoring & 0.30 \\
Phonetic values & 0.03 \\
Language identification & 0.06 \\
Strict full decipherment & 0.10 \\
\bottomrule
\end{tabular}
\caption{Heuristic maturity estimates. These are not probabilities of final success.}
\label{tab:decipherment-distance}
\end{table}

The honest answer is therefore: we may be roughly halfway toward a serious structural decipherment testbench, but only around a tenth of the way toward strict full decipherment. The fastest safe path is not to guess a language now, but to build a falsifiable chain:

\begin{center}
\texttt{ANCHORS -> COMPONENTS -> PROPER NAMES/TITLES -> PHONETIC TESTS}
\end{center}

\subsection{Next Research Decision}

The next milestone should expand the solver from the eleven anchor units to their nearest minimal-pair neighbors. This directly attacks the phonetic problem: if \texttt{240-176}, \texttt{240-482}, and related alternants preserve frame identity while substituting only one component, then the project can begin testing whether the contrast behaves like a title, stem, suffix, numeral, allograph, or phonetic alternation.

\subsection{Outputs}

\begin{itemize}
  \item \texttt{outputs/constraint\_solver\_summary.csv}
  \item \texttt{outputs/constrained\_reading\_candidates.csv}
  \item \texttt{outputs/constraint\_violations.csv}
  \item \texttt{outputs/morpheme\_slot\_assignments.csv}
  \item \texttt{outputs/reconstructed\_clause\_frames.csv}
  \item \texttt{outputs/decipherment\_progress\_estimate.csv}
  \item \texttt{outputs/constraint\_solver\_model.tex}
\end{itemize}
